\documentclass{hswbeamer}

\title[EiP]{Einführung in die Programmierung}
\subtitle[WS 2014/15]{Wintersemester 2014/15}
\author[\Copyright{Uwe L\"{a}mmel, Thomas Jonitz}]{Prof. Dr.-Ing Uwe L\"{a}mmel \and Thomas Jonitz (B.Sc.)}
\date{Wismar, \today}
%\logo{\includegraphics[height=1cm]{fischer.png}}
\titlegraphic{\includegraphics[scale=0.05]{img/hsw-wings.jpg}}

\bibliography{literatur}
\makeindex
\begin{document}



\begin{frame}
    \begin{quote}
        \glqq Saying Java is nice, because it works on all OSes ist like saying anal sex is nice, because it works on all genders.\grqq -- Icematix (c-plusplus.de)
    \end{quote}
\end{frame}

\begin{frame}
    \maketitle
\end{frame}

\begin{frame}{Inhaltsverzeichnis}{Agenda}
    \begin{columns}
        \begin{column}{0.4\textwidth}
            \includegraphics[width=0.9\columnwidth]{img/java-tm.jpg}
        \end{column}
        \begin{column}{0.6\textwidth}
            \tableofcontents%[hideothersections]
        \end{column}
    \end{columns}
\end{frame}

\section{Einführung}\label{sec:0-einfuehrung}
\begin{frame}{Einführung}
\end{frame}


\begin{frame}{Ablauf der Lehrveranstaltung}
\end{frame}

\begin{frame}{Motto}
\end{frame}

\begin{frame}{Problem}
\end{frame}

\begin{frame}{Literaturempfehlung}
\end{frame}
\section{Klassendefinition und Objekte}\label{sec:1-objekte}
\begin{frame}{}
\begin{center}
\Huge Klassendefinition und Objekte
\end{center}
\end{frame}

\section{Anweisungen und Datentypen}
\begin{frame}{Sample}
\end{frame}

\section{Datensammlungen}
\begin{frame}{Sample}
\end{frame}

\section{Container-Klassen}
\begin{frame}{Sample}
\end{frame}

\section{Algorithmen}
\begin{frame}{Sample}
\end{frame}

\section{Vererbung}
\begin{frame}{Sample}
\end{frame}

\section{Statische Methoden}
\begin{frame}{Sample}
\end{frame}

\section{Zusammenfassung}
\begin{frame}{Sample}
\end{frame}



\end{document}
